% questions.tex —— 2013 年·陕西卷·理科 · 题目内容片段（无答案）
% 由 tools/generate_projects.py 依来源总账自 2013/2013_数学真题_陕西理_试题.md 生成，勿手改（改源 MD 后重跑）


\begin{gknotice}
\textbf{来源说明}：本题面依据固定源 PDF 的两页 A3 题面逐题转录，原卷印刷页码为 837—838。第 837 面为选择题 1—10、填空题 11—17、解答题 18—19，第 838 面为解答题 20—23。第 15—17 题分别对应不等式选讲、几何证明选讲、坐标系与参数方程，按原卷分别编号并全部收录。

\textbf{解析归属}：固定源不附答案或解析；配套答案与解析均由本库按本题面独立复算。

\textbf{图形说明}：第 5 题矩形内两个扇形、第 9 题内接矩形、第 12 题三视图、第 16 题圆内相交弦、第 17 题圆的参数方程示意、第 20 题四棱柱，均由项目内 TikZ 矢量源重绘。
\end{gknotice}

\gkpart{一、选择题}

\begin{question}
\gkti{1}设全集为 $\mathbf R$，函数 $f(x)=\sqrt{1-x^2}$ 的定义域为 $M$，则 $\complement_{\mathbf R}M$ 为\mbox{（　　）}

（A）$[-1,1]$；（B）$(-1,1)$；（C）$(-\infty,-1]\cup[1,+\infty)$；（D）$(-\infty,-1)\cup(1,+\infty)$。

\end{question}
\begin{question}
\gkti{2}根据下列算法语句，当输入 $x$ 为 $60$ 时，输出 $y$ 的值为\mbox{（　　）}

\nopagebreak[4]
\[\begin{array}{|l|}
\hline
\text{\text{输入} }x;\\
\text{If }x\leqslant 50\text{ Then}\\
\qquad y=0.5*x\\
\text{Else}\\
\qquad y=25+0.6*(x-50)\\
\text{End If}\\
\text{\text{输出} }y.\\
\hline
\end{array}\]

（A）$25$；（B）$30$；（C）$31$；（D）$61$。

\end{question}
\begin{question}
\gkti{3}设 $\boldsymbol a$，$\boldsymbol b$ 为向量，则“$|\boldsymbol a\cdot\boldsymbol b|=|\boldsymbol a||\boldsymbol b|$”是“$\boldsymbol a\parallel\boldsymbol b$”的\mbox{（　　）}

（A）充分不必要条件；（B）必要不充分条件；（C）充分必要条件；（D）既不充分也不必要条件。

\end{question}
\begin{question}
\gkti{4}某单位有 $840$ 名职工，现采用系统抽样方法，抽取 $42$ 人做问卷调查，将 $840$ 人按 $1,2,\cdots,840$ 随机编号，则抽取的 $42$ 人中，编号落入区间 $[481,720]$ 的人数为\mbox{（　　）}

（A）$11$；（B）$12$；（C）$13$；（D）$14$。

\Needspace{0.4\textheight}
\end{question}
\begin{question}
\gkti{5}如图，在矩形区域 $ABCD$ 的 $A$，$C$ 两点处各有一个通信基站，假设其信号覆盖范围分别是扇形区域 $ADE$ 和扇形区域 $CBF$（该矩形区域内无其他信号来源，基站工作正常）。若在该矩形区域内随机地选一地点，则该地点无信号的概率是\mbox{（　　）}
\nopagebreak[4]
\begin{center}% ============================================================
% f2013_shaanxi_science_q05_sectors.tex
% 2013 陕西卷（理工类）第 5 题 矩形内的两个扇形（TikZ 矢量重绘）
% 矩形 ABCD 中 AD=1、AB=2；扇形 ADE 与 CBF 的半径均为 1
% ============================================================
\begin{tikzpicture}[
  x=2.2cm, y=2.2cm, line width=0.5pt, font=\footnotesize, >=Stealth
]

  % ---------- 矩形与扇形弧 ----------
  \draw (0,0) rectangle (2,1);
  \draw (0,1) arc (90:0:1);      % 以 A 为圆心，D 到 E
  \draw (1,1) arc (180:270:1);   % 以 C 为圆心，F 到 B

  % ---------- 顶点 ----------
  \node[below left]  at (0,0) {$A$};
  \node[below right] at (2,0) {$B$};
  \node[above right] at (2,1) {$C$};
  \node[above left]  at (0,1) {$D$};
  \node[above]       at (1,1) {$F$};
  \node[above right] at (1,0) {$E$};

  % ---------- 尺寸标注 ----------
  \draw[|<->|] (-0.22,0) -- (-0.22,1) node[midway, left] {$1$};
  \foreach \y in {0,1} \draw (-0.32,\y) -- (-0.12,\y);
  \draw[|<->|] (0,-0.22) -- (2,-0.22) node[midway, below] {$2$};
  \foreach \x in {0,2} \draw (\x,-0.32) -- (\x,-0.12);

\end{tikzpicture}
\end{center}
\nopagebreak[4]
（A）$1-\dfrac{\pi}{4}$；（B）$\dfrac{\pi}{2}-1$；（C）$2-\dfrac{\pi}{2}$；（D）$\dfrac{\pi}{4}$。

\end{question}
\begin{question}
\gkti{6}设 $z_1$，$z_2$ 是复数，则下列命题中的假命题是\mbox{（　　）}

（A）若 $|z_1-z_2|=0$，则 $z_1=z_2$；（B）若 $z_1=\bar z_2$，则 $\bar z_1=z_2$；（C）若 $|z_1|=|z_2|$，则 $z_1\cdot\bar z_1=z_2\cdot\bar z_2$；（D）若 $|z_1|=|z_2|$，则 $z_1^2=z_2^2$。

\end{question}
\begin{question}
\gkti{7}设 $\triangle ABC$ 的内角 $A$，$B$，$C$ 所对的边分别为 $a$，$b$，$c$，若 $b\cos C+c\cos B=a\sin A$，则 $\triangle ABC$ 的形状为\mbox{（　　）}

（A）锐角三角形；（B）直角三角形；（C）钝角三角形；（D）不确定。

\end{question}
\begin{question}
\gkti{8}设函数
\nopagebreak[4]
\[f(x)=\begin{cases}
\left(x-\dfrac1x\right)^6,&x<0,\\[6pt]
-\sqrt x,&x\geqslant0,
\end{cases}\]
则当 $x>0$ 时，$f[f(x)]$ 表达式的展开式中常数项为\mbox{（　　）}

（A）$-20$；（B）$20$；（C）$-15$；（D）$15$。

\Needspace{0.4\textheight}
\end{question}
\begin{question}
\gkti{9}在如图所示的锐角三角形空地中，欲建一个面积不小于 $300\ \mathrm{m}^2$ 的内接矩形花园（阴影部分），则其边长 $x$（单位 m）的取值范围是\mbox{（　　）}
\nopagebreak[4]
\begin{center}% ============================================================
% f2013_shaanxi_science_q09_triangle.tex
% 2013 陕西卷（理工类）第 9 题 锐角三角形空地与内接矩形（TikZ 矢量重绘）
% 三角形底边 40 m、高 40 m；内接矩形上边长为 x
% ============================================================
\begin{tikzpicture}[
  x=0.13cm, y=0.13cm, line width=0.5pt, font=\footnotesize, >=Stealth
]

  % ---------- 内接矩形（阴影） ----------
  \fill[gray!18] (7.5,0) rectangle (32.5,15);

  % ---------- 三角形空地 ----------
  \draw (0,0) -- (40,0) -- (20,40) -- cycle;

  % ---------- 内接矩形边界 ----------
  \draw (7.5,0) -- (7.5,15) -- (32.5,15) -- (32.5,0);

  % ---------- 标注 ----------
  \node[above] at (20,15.6) {$x$};

  \draw[|<->|] (43,0) -- (43,40) node[midway, right] {$40\ \mathrm{m}$};
  \foreach \y in {0,40} \draw (41.6,\y) -- (44.4,\y);
  \draw[|<->|] (0,-4.2) -- (40,-4.2) node[midway, below] {$40\ \mathrm{m}$};
  \foreach \x in {0,40} \draw (\x,-2.6) -- (\x,-5.8);

\end{tikzpicture}
\end{center}
\nopagebreak[4]
（A）$[15,20]$；（B）$[12,25]$；（C）$[10,30]$；（D）$[20,30]$。

\end{question}
\begin{question}
\gkti{10}设 $[x]$ 表示不大于 $x$ 的最大整数，则对任意实数 $x$，$y$，有\mbox{（　　）}

（A）$[-x]=-[x]$；（B）$[2x]=2[x]$；（C）$[x+y]\leqslant[x]+[y]$；（D）$[x-y]\leqslant[x]-[y]$。

\end{question}
\begingroup
\renewcommand{\gkpart}[1]{\Needspace{0.2\textheight}\subsection*{#1}\nopagebreak}
\gkpart{二、填空题}
\endgroup

\begin{question}
\gkti{11}双曲线 $\dfrac{x^2}{16}-\dfrac{y^2}{m}=1$ 的离心率为 $\dfrac54$，则 $m$ 等于\gkblank{}。

\Needspace{0.4\textheight}
\end{question}
\begin{question}
\gkti{12}某几何体的三视图如图所示，则其体积为\gkblank{}。
\nopagebreak[4]
\begin{center}% ============================================================
% f2013_shaanxi_science_q12_views.tex
% 2013 陕西卷（理工类）第 12 题 三视图（TikZ 矢量重绘）
% 主视图：直角边 1、2 的直角三角形
% 左视图：底边 2、高 2 的等腰三角形
% 俯视图：半径 1 的半圆
% 还原为底面半径 1、高 2 的半圆锥，体积 π/3
% ============================================================
\begin{tikzpicture}[
  x=1cm, y=1cm, line width=0.5pt, font=\footnotesize, >=Stealth
]

  % ---------- 主视图 ----------
  \draw (0,0) -- (1,0) -- (0,2) -- cycle;
  \draw[|<->|] (-0.35,0) -- (-0.35,2) node[midway, left] {$2$};
  \foreach \y in {0,2} \draw (-0.50,\y) -- (-0.20,\y);
  \draw[|<->|] (0,-0.35) -- (1,-0.35) node[midway, below] {$1$};
  \foreach \x in {0,1} \draw (\x,-0.50) -- (\x,-0.20);
  \node at (0.5,-1.00) {主视图};

  % ---------- 左视图 ----------
  \begin{scope}[xshift=2.6cm]
    \draw (0,0) -- (2,0) -- (1,2) -- cycle;
    \draw[|<->|] (0,-0.35) -- (1,-0.35) node[midway, below] {$1$};
    \draw[|<->|] (1,-0.35) -- (2,-0.35) node[midway, below] {$1$};
    \foreach \x in {0,1,2} \draw (\x,-0.50) -- (\x,-0.20);
    \node at (1,-1.00) {左视图};
  \end{scope}

  % ---------- 俯视图 ----------
  \draw (0.5,-3.20) -- (0.5,-5.20);
  \draw (0.5,-3.20) arc (90:-90:1);
  \node at (1.20,-5.65) {俯视图};

\end{tikzpicture}
\end{center}
\nopagebreak[4]
\end{question}
\begin{question}
\gkti{13}若点 $(x,y)$ 位于曲线 $y=|x-1|$ 与 $y=2$ 所围成的封闭区域，则 $2x-y$ 的最小值为\gkblank{}。

\end{question}
\begin{question}
\gkti{14}观察下列等式：
\nopagebreak[4]
\[\begin{aligned}
1^2&=1,\\
1^2-2^2&=-3,\\
1^2-2^2+3^2&=6,\\
1^2-2^2+3^2-4^2&=-10,\\
&\cdots\cdots
\end{aligned}\]
照此规律，第 $n$ 个等式可为\gkblank{}。

\end{question}
\begin{question}
\gkti{15}已知 $a$，$b$，$m$，$n$ 均为正数，且 $a+b=1$，$mn=2$，则 $(am+bn)(bm+an)$ 的最小值为\gkblank{}。

\Needspace{0.4\textheight}
\end{question}
\begin{question}
\gkti{16}如图，弦 $AB$ 与 $CD$ 相交于 $\odot O$ 内一点 $E$，过 $E$ 作 $BC$ 的平行线与 $AD$ 的延长线相交于点 $P$。已知 $PD=2DA=2$，则 $PE=$\gkblank{}。
\nopagebreak[4]
\begin{center}% ============================================================
% f2013_shaanxi_science_q16_geometry.tex
% 2013 陕西卷（理工类）第 16 题 圆内相交弦与平行线（TikZ 矢量重绘）
% 弦 AB、CD 交于 E；过 E 作 BC 的平行线与 AD 的延长线交于 P
% PD=2DA=2，按此比例作图
% ============================================================
\begin{tikzpicture}[
  x=1cm, y=1cm, line width=0.5pt, font=\footnotesize, >=Stealth
]

  % ---------- 圆 ----------
  \coordinate (O) at (5.50,1.739);
  \draw (O) circle (2.058);

  % ---------- 点 ----------
  \coordinate (P) at (0,0);
  \coordinate (D) at (4.40,0);
  \coordinate (A) at (6.60,0);
  \coordinate (B) at (3.576,2.469);
  \coordinate (C) at (6.984,3.165);
  \coordinate (E) at (5.28,1.078);

  % ---------- 线段 ----------
  \draw (P) -- (A);
  \draw (A) -- (B);
  \draw (C) -- (D);
  \draw (B) -- (C);
  \draw (P) -- (E);

  % ---------- 圆心 ----------
  \fill (O) circle (0.7pt);

  % ---------- 标注 ----------
  \node[below left]  at (P) {$P$};
  \node[below]       at (D) {$D$};
  \node[below right] at (A) {$A$};
  \node[left]        at (B) {$B$};
  \node[right]       at (C) {$C$};
  \node[below right] at (E) {$E$};
  \node[above right] at (O) {$O$};

\end{tikzpicture}
\end{center}
\nopagebreak[4]
\Needspace{0.4\textheight}
\end{question}
\begin{question}
\gkti{17}如图，以过原点的直线的倾斜角 $\theta$ 为参数，则圆 $x^2+y^2-x=0$ 的参数方程为\gkblank{}。
\nopagebreak[4]
\begin{center}% ============================================================
% f2013_shaanxi_science_q17_circle.tex
% 2013 陕西卷（理工类）第 17 题 圆的参数方程示意（TikZ 矢量重绘）
% 圆 x^2+y^2-x=0，过原点的直线的倾斜角为 θ，交圆于 P
% ============================================================
\begin{tikzpicture}[
  x=3cm, y=3cm, line width=0.5pt, font=\footnotesize, >=Stealth
]

  % ---------- 坐标轴 ----------
  \draw[->] (-0.16,0) -- (1.16,0) node[right] {$x$};
  \draw[->] (0,-0.16) -- (0,0.96) node[above] {$y$};

  % ---------- 圆 ----------
  \draw (0.5,0) circle (0.5);

  % ---------- 过原点的直线与交点 P ----------
  \coordinate (P) at (0.671,0.470);
  \draw (0,0) -- (P);
  \fill (P) circle (0.8pt);

  % ---------- 角 θ ----------
  \draw (0.160,0) arc (0:35:0.160);
  \node at (0.250,0.072) {$\theta$};

  % ---------- 标注 ----------
  \node[below left]  at (0,0) {$O$};
  \node[above right] at (P) {$P$};

\end{tikzpicture}
\end{center}
\nopagebreak[4]
\end{question}
\begingroup
\renewcommand{\gkpart}[1]{\Needspace{0.2\textheight}\subsection*{#1}\nopagebreak}
\gkpart{三、解答题}
\endgroup

\begin{question}
\gkti{18}已知向量 $\boldsymbol a=\left(\cos x,-\dfrac12\right)$，$\boldsymbol b=\left(\sqrt3\sin x,\cos2x\right)$，$x\in\mathbf R$，设函数 $f(x)=\boldsymbol a\cdot\boldsymbol b$。

（1） 求 $f(x)$ 的最小正周期。

（2） 求 $f(x)$ 在 $\left[0,\dfrac\pi2\right]$ 上的最大值和最小值。

\end{question}
\begin{question}
\gkti{19}设 $\{a_n\}$ 是公比为 $q$ 的等比数列。

（1） 推导 $\{a_n\}$ 的前 $n$ 项和公式；

（2） 设 $q\neq1$，证明数列 $\{a_n+1\}$ 不是等比数列。

\Needspace{0.4\textheight}
\end{question}
\begin{question}
\gkti{20}如图，四棱柱 $ABCD-A_1B_1C_1D_1$ 的底面 $ABCD$ 是正方形，$O$ 为底面中心，$A_1O\perp$ 平面 $ABCD$，$AB=AA_1=\sqrt2$。

（1） 证明：$A_1C\perp$ 平面 $BB_1D_1D$；

（2） 求平面 $OCB_1$ 与平面 $BB_1D_1D$ 的夹角 $\theta$ 的大小。
\nopagebreak[4]
\begin{center}% ============================================================
% f2013_shaanxi_science_q20_prism.tex
% 2013 陕西卷（理工类）第 20 题 四棱柱（TikZ 矢量重绘）
% 底面 ABCD 为正方形，O 为底面中心，A1O ⊥ 平面 ABCD
% 虚线：被遮挡的底边、侧棱、A1O 以及 A1C、BD
% ============================================================
\begin{tikzpicture}[
  x=1cm, y=1cm, line width=0.5pt, font=\footnotesize, >=Stealth
]

  % ---------- 顶点 ----------
  \coordinate (A)  at (0,0);
  \coordinate (B)  at (3,0);
  \coordinate (C)  at (4.5,1.1);
  \coordinate (D)  at (1.5,1.1);
  \coordinate (O)  at (2.25,0.55);
  \coordinate (A1) at (2.25,3.6);
  \coordinate (B1) at (5.25,3.6);
  \coordinate (C1) at (6.75,4.7);
  \coordinate (D1) at (3.75,4.7);

  % ---------- 底面 ----------
  \draw (A) -- (B) -- (C);
  \draw[dashed] (C) -- (D) -- (A);

  % ---------- 侧棱 ----------
  \draw (A) -- (A1);
  \draw (B) -- (B1);
  \draw (C) -- (C1);
  \draw[dashed] (D) -- (D1);

  % ---------- 上底面 ----------
  \draw (A1) -- (B1) -- (C1) -- (D1) -- cycle;

  % ---------- 高与辅助线 ----------
  \draw[dashed] (A1) -- (O);
  \draw[dashed] (A1) -- (C);
  \draw[dashed] (B) -- (D);

  % ---------- 标注 ----------
  \node[below left]  at (A)  {$A$};
  \node[below right] at (B)  {$B$};
  \node[right]       at (C)  {$C$};
  \node[left]        at (D)  {$D$};
  \node[left]        at (A1) {$A_1$};
  \node[above]       at (B1) {$B_1$};
  \node[above right] at (C1) {$C_1$};
  \node[above]       at (D1) {$D_1$};
  \node[below right] at (O)  {$O$};

\end{tikzpicture}
\end{center}
\nopagebreak[4]
\end{question}
\begin{question}
\gkti{21}在一场娱乐晚会上，有 $5$ 位民间歌手（$1$ 至 $5$ 号）登台演唱，由现场数百名观众投票选出最受欢迎歌手。各位观众须彼此独立地在选票上选 $3$ 名歌手，其中观众甲是 $1$ 号歌手的歌迷，他必选 $1$ 号，不选 $2$ 号，另在 $3$ 至 $5$ 号中随机选 $2$ 名。观众乙和丙对 $5$ 位歌手的演唱没有偏爱，因此在 $1$ 至 $5$ 号中随机选 $3$ 名歌手。

（1） 求观众甲选中 $3$ 号歌手且观众乙未选中 $3$ 号歌手的概率；

（2） $X$ 表示 $3$ 号歌手得到观众甲、乙、丙的票数之和，求 $X$ 的分布列和数学期望。

\end{question}
\begin{question}
\gkti{22}已知动圆过定点 $A(4,0)$，且在 $y$ 轴上截得的弦 $MN$ 的长为 $8$。

（1） 求动圆圆心的轨迹 $C$ 的方程；

（2） 已知点 $B(-1,0)$，设不垂直于 $x$ 轴的直线 $l$ 与轨迹 $C$ 交于不同的两点 $P$，$Q$，若 $x$ 轴是 $\angle PBQ$ 的角平分线，证明直线 $l$ 过定点。

\end{question}
\begin{question}
\gkti{23}已知函数 $f(x)=\mathrm e^x$，$x\in\mathbf R$。

（1） 若直线 $y=kx+1$ 与 $f(x)$ 的反函数的图象相切，求实数 $k$ 的值；

（2） 设 $x>0$，讨论曲线 $y=f(x)$ 与曲线 $y=mx^2$（$m>0$）公共点的个数；

（3） 设 $a<b$，比较 $\dfrac{f(a)+f(b)}2$ 与 $\dfrac{f(b)-f(a)}{b-a}$ 的大小，并说明理由。

\end{question}